% !TEX root = ../main.tex
\chapter{Referência de diretivas}
\label{ap:diretivas}

\section{Diretivas de configuração}

Uma por linha, no topo do \file{.cmm}, sem ponto e vírgula. A coluna do padrão indica o valor assumido pela cadeia de ferramentas quando a diretiva é omitida; o Hub de Processadores da AURORA escreve todas explicitamente.

\begin{longtable}{lllp{0.42\linewidth}}
\toprule
Diretiva & Argumento & Padrão & Efeito \\
\midrule
\endhead
\code{\#PRNAME} & nome & & Nome do processador; deve casar com o nome no projeto e nos artefatos. \\
\code{\#NUBITS} & inteiro & 23 & Largura da palavra da ULA e dos inteiros, em complemento de dois; deve valer a soma de mantissa, expoente e um. \\
\code{\#NBMANT} & inteiro & 16 & Bits de mantissa do ponto flutuante próprio. \\
\code{\#NBEXPO} & inteiro & 6 & Bits de expoente, em complemento de dois. \\
\code{\#NDSTAC} & inteiro & 10 & Profundidade da pilha de dados. \\
\code{\#SDEPTH} & inteiro & 10 & Profundidade da pilha de sub-rotinas. \\
\code{\#NUIOIN} & inteiro & 1 & Número de portas de entrada; \code{in(p)} exige \texttt{p} menor que esse valor. \\
\code{\#NUIOOU} & inteiro & 1 & Número de portas de saída; idem para \texttt{out}. \\
\code{\#NUGAIN} & potência de 2 & 64 & Divisor fixo da função \code{norm()}. \\
\code{\#FFTSIZ} & inteiro & 8 & Bits invertidos no endereçamento bit-reverso \code{x[k)}, para FFT de $2^n$ pontos. \\
\bottomrule
\end{longtable}

Com uma única porta a ligação é direta; com mais de uma, o hardware instancia automaticamente um decodificador de endereços.

\section{Diretivas comportamentais}

Aparecem no corpo do programa, marcando pontos do código. A \code{\#PRACA} marca o ponto para onde o processador desvia quando o pino de interrupção pulsa. A \code{\#TOAQUI} marca um endereço do programa cujo alcance faz pulsar o pino \code{cheguei}; é também o mecanismo que a AURORA usa para detectar o fim do programa em testes e simulações.

\section{Pré-processador}

A construção \code{\#define NOME corpo} define uma constante simbólica, apenas na forma de objeto, sem argumentos, com o limite de 256 por programa. Não há \code{\#include} nem \code{\#ifdef} no fluxo \cmm{}.
